\appendix

\section{Appendix}

\subsection{Hiding Monad on Heaplets}
\label{adx:heaplet-hiding-mon}

Recall the action of $\dhideMon$ on objects:
\[\begin{array}{r@{\,}@{\,}l@{\,}l@{\,}}
    \dhideMon          &:&\funC\dinitC\setC \to \funC\dinitC\setC \\
    (\dhideMon\vA)(\rho : \vw^- \injto \vw^+)  &=& \Coend{\vw^- \injto \vw' \in \commaC{\vw^-}\vU} 
                                                             \vA(\vi^{\vw' \inclto \vw' \oplus \vw^+ \ominus \rho})
\end{array}\]
The action on morphism $\varepsilon = 
\begin{tikzcd}
    {\vw_1^+} \arrow[rr, "{\vU_2\varepsilon}"] && {\vw_2^+} \\
    {\vw_1^-} \arrow[u, "{\rho}"] \arrow[rr, "{\vU_1\varepsilon}"'] && {\vw_2^-} \arrow[u, "{\sigma}"'] \\
\end{tikzcd}, \vh_\varepsilon : \dheapF(\vw_2^- \ominus \vU_1\varepsilon, \vw_2^+)$ is given as:

\[\begin{array}{r@{\,}@{\,}l@{\,}l@{\,}}
    \dhideMon\vA(\varepsilon) &:& \dhideMon\vA\rho \to \dhideMon\vA\sigma \\
    \quot{\vw_1^- \xrightarrow{\rho'} \vw'}\va & \mapsto&
        \quot{\vw_2^- \xrightarrow{(\vU_1\varepsilon)^*\rho'} \rho' \oplus_{\vw_1^-} \vU_1\varepsilon}\vA(\varepsilon')(\va)
\end{array}\]

Where $(\vU_1\varepsilon)^*\rho'$ comes from:
\begin{tikzcd}
    {\vw_1^-} \arrow[d, "\rho'"'] \arrow[rr, "{\vU_1\varepsilon}"] && {\vw_2^-} \arrow[d, "{(\vU_1\varepsilon)^*\rho'}"] \\
    {\vw'} \arrow[rr, "{(\rho')^*\vU_1\varepsilon}"'] && {\rho' \oplus_{\vw_1^-} \vU_1\varepsilon}
\end{tikzcd}

and $\varepsilon'$ is:
\begin{tikzcd}
    {\vw' \oplus (\vw_1^+ \ominus \rho)} \arrow[rrr, "{[(\rho')^*\vU_1\varepsilon, \psi \comp \vU_2\varepsilon \comp \compl\rho]}"] &&& 
        {\rho' \oplus_{\vw_1^-} \vU_1\varepsilon \oplus \vw_2^+ \ominus \sigma} \\
    {\vw'} \arrow[u, hook] \arrow[rrr, "{(\rho')^*\vU_1\varepsilon}"] &&& 
        {\rho' \oplus_{\vw_1^-} \vU_1\varepsilon} \arrow[u, hook]
\end{tikzcd}, $\dheapF(\iso, \psi)(\eta_\varepsilon)$

where $\psi$ is given by: 
\[ \vw_2^+ \xrightarrow\iso \vw_1^- \oplus (\vw_2^- \ominus \vU_1\varepsilon) \oplus (\vw_2^+ \ominus \sigma)
           \inclto \rho' \oplus_{\vw_1^-} \vU_1\varepsilon \oplus \vw_2^+ \ominus \sigma
\]

We will often consider the action of $\dhideMon\vA(\varepsilon)$ as two parts: 
\begin{itemize}
    \item An action between elements of the equivalence classes (does not satisfy functoriality): 
        \[ \begin{array}{r@{\,}c@{\,}l}
            \dhideMonSk\vA(\varepsilon : \rho \to \sigma) & : & \Coprod{\vw_1^- \injto \vw'} \vA(\inl^{\vw' \to \vw' \oplus (\vw_1^+ \ominus \rho)}) \to
                                                                \Coprod{\vw_2^- \injto \vw'} \vA(\inl^{\vw' \to \vw' \oplus (\vw_2^+ \ominus \sigma)}) \\
            \dhideMonSk\vA(\varepsilon)(\inj{\rho'}\va) &  =  & \inj{\vU_1\varepsilon*\rho'}\vA(\varepsilon')(\va)
        \end{array}\]
    \item A quotienting step given by the coend to make it functorial: 
            \[ \dhideMon\vA(\varepsilon)(\quot{\rho'}\va) = [\underline{\mathbb{P}}\vA(\varepsilon)(\inj{\rho'}\va)] \]
\end{itemize}

% \subsection{Heaplet Compatibility}
% \label{adx:heaplet-compatibility}

% We say that $\tup{\vw \xrightarrow\rho \vw^- \xrightarrow{\rho^+} \vw^+}{\vh' : \dheapF(\vw^-,\vw^+)}$ is compatible with:
% \[ \tup{{\ubar\vw}^- \xrightarrow{\sigma^+} {\ubar\vw}^+ \xleftarrow\sigma \vw^+}{\ubar\vh : \dheapF(\ubar\vw^-, \ubar\vw^+)} \]
% when the following conditions are met:
% \begin{itemize}
%     \item For all their shared locations, $\vh'$ and $\ubar\vh$ store the same values, that is, we have:
%             \[ \dheapF(\nu_*, \sigma)(\vh') = \dheapF(\ext\nu, 1_{\ubar\vw^+})(\ubar\vh) \]

%         Where $\nu_*, \ext\nu$ come from the following pullback diagram:
%         \begin{center}
%             \begin{tikzcd}
%                 {\vv} \arrow[d, "{\nu_*}"'] \arrow[rr, "{\ext\nu}"] && {\ubar\vw^-} \arrow[d, "{\sigma^+}"] \\
%                 {\vw^-} \arrow[r, "{\rho^+}"] & {\vw^+} \arrow[r, "{\sigma}"] & {\ubar\vw^+} \\
%             \end{tikzcd}
%         \end{center}
%     \item The part of heap $\ubar\vh$ that is not shared with $\vh'$ does not have references to locations $\vw^- \ominus \rho$ (representing the newly allocated locations):

%           There exists a $\eta_{\varepsilon_1} : \dheapF(\ubar\vw^- \ominus \nu^*, \ubar\vw^+ \ominus \tau)$ where $\tau$ is:
%             \[ \vw^- \ominus \rho \xrightarrow{\compl\rho} \vw^- \xrightarrow{\rho^+} \vw^+ \xrightarrow\sigma \ubar\vw^+ \]
%           such that $\dheapF(1, \compl\tau)(\eta_{\varepsilon_1}) = \dheapF(\compl{\ext\nu}, 1)(\ubar\vh)$.
% \end{itemize}

% % The above data give us the following morphisms in $\dinitC$:
% % \[\begin{array}{lcr}
% %     \varepsilon_1 : 1_\vv \to \sigma^+ = \triple{\ext\nu}{\sigma^+ \comp \ext\nu}{\dheapF(1, \compl\tau)(\eta_{\varepsilon_1})} & \text{ and } &
% %     \varepsilon_2 : 1_\vv \to \rho^+ = \triple{\nu_*}{\rho^+ \comp \nu_*}{\dheapF(\compl{\nu_*}, 1)(\vh')}
% % \end{array}\]

% The above data give us the following heaplets:
% \[\begin{array}{lcr}
%     \eta_{\varepsilon_1} : \dheapF(\ubar\vw^- \ominus \nu^*, \ubar\vw^+ \ominus \tau) & \text{ and } &
%     \eta_{\varepsilon_2} = \dheapF(\compl{\nu_*}, 1)(\vh') : \dheapF(\vw^- \ominus \nu_*, \vw^+)
% \end{array}\]

\subsection{End related lemmas}
\label{adx:end-lemmas}

Unless otherwise specified let $\vA,\vB,\vG : \mcatC \to \setC$, $\vF : \mcatC \to \mcatD$  and the following morphisms in $\commaC\vD\vF$:
\begin{center}
\begin{tikzcd}
    {\vF\vC}
        \arrow[rr, "{\vF\vh}"]
    &&
    {\vF\vC'}
    \\
    &
    {\vD}
        \arrow[ul, "\vf"]
        \arrow[ur, "\vg"']
\end{tikzcd}
\end{center}

\begin{lemma}\label{adx:coslice-ends-proof}
    \[ \End{\vC \in \mcatC} \Exp{\vG\vC}{\homF\mcatD\vD{\vF\vC}} \iso \End{\vf : \vD \to \vF\vC \in \commaC\vD\vF} \vG\vC \]
\end{lemma}
\begin{proof}
    The morphism $\End{\vC \in \mcatC} \Exp{\vG\vC}{\homF\mcatD\vD{\vF\vC}} \dashrightarrow \End{\vf : \vD \to \vF\vC \in \commaC\vD\vF} \vG\vC$ is from the following wedge:
    \begin{center}
        \begin{tikzcd}
            & {\End{\vC \in \mcatC} \Exp{\vG\vC}{\homF\mcatD\vD{\vF\vC}}}
                \arrow[dl, "\proj\vC"']
                \arrow[dr, "\proj{\vC'}"]
        \\
            {\Exp{\vG\vC}{\homF\mcatD\vD{\vF\vC}}}
                \arrow[dd, "{\dash(\vf)}"]
                \arrow[dr, "{\vG\vh \comp \dash}"]
            &&
            {\Exp{\vG{\vC'}}{\homF\mcatD\vD{\vF{\vC'}}}}
                \arrow[dd, "{\dash(\vg)}"]
                \arrow[dl, "{\Lam\vk \dash(\vF\vh \comp \vk)}"']
        \\
            &
                {\Exp{\vG{\vC'}}{\homF\mcatD\vD{\vF\vC}}}
                    \arrow[dr, "{\dash(\vf)}"]
            &
        \\
            {\vG\vC}
                \arrow[rr, "{\vG\vh}"]
            &&
            {\vG{\vC'}}
        \end{tikzcd}
    \end{center}
    The morphism $\End{\vf : \vD \to \vF\vC \in \commaC\vD\vF} \vG\vC \dashrightarrow \End{\vC \in \mcatC} \Exp{\vG\vC}{\homF\mcatD\vD{\vF\vC}} $ is from the following wedge:
    \begin{center}
        \begin{tikzcd}
            &
                {\End{\vf : \vD \to \vF\vC \in \commaC\vD\vF} \vG\vC}
                    \arrow[dl, "{\Lam\vf \proj\vf(\dash)}"']
                    \arrow[dr, "{\Lam\vg \proj\vg(\dash)}"]
                    \arrow[dd, "{\Lam\vk \proj{\vF\vh \comp \vk}(\dash)}"]
        \\
            {\Exp{\vG\vC}{\homF\mcatD\vD{\vF\vC}}}
                \arrow[dr, "{\vG\vh \comp \dash}"]
            &&
            {\Exp{\vG{\vC'}}{\homF\mcatD\vD{\vF{\vC'}}}}
                \arrow[dl, "{\Lam\vk \dash(\vF\vh \comp \vk)}"]
        \\
            &
                {\Exp{\vG{\vC'}}{\homF\mcatD\vD{\vF\vC}}}
        \end{tikzcd}
    \end{center}
    It follows directly that these two induced morphisms are inverses, and thus the two ends are isomorphic.
\end{proof}

\begin{lemma}\label{adx:coslice-end-exponent-proof}
    \[ \End{\vf : \vD \to \vF\vC \in \commaC\vD\vF} \Exp\vB\vA\vC \iso \End{\vf : \vD \to \vF\vC \in \commaC\vD\vF} \Exp{\vB\vC}{\vA\vC} \]
\end{lemma}
\begin{proof}

The morphism  $\End{\vf : \vD \to \vF\vC \in \commaC\vD\vF} \Exp\vB\vA\vC \dashrightarrow \End{\vf : \vD \to \vF\vC \in \commaC\vD\vF} \Exp{\vB\vC}{\vA\vC}$ is from the following wedge:
\begin{center}
    \begin{tikzcd}
    & {\End{\vf : \vD \to \vF\vC \in \commaC\vD\vF} \Exp\vB\vA\vC}
        \arrow[dl, "{\proj\vf}"']
        \arrow[dr, "{\proj\vg}"]
    \\
    {\Exp\vB\vA\vC}
        \arrow[d, "{\Lam\va (\dash)_\vC\tup{1_\vC}\va}"]
        \arrow[rr, "{(\Lam{\tup\vk\va} (\dash)_{\hat\vC}\tup{\vk \comp \vh}\va)_{\hat\vC \in \mcatC}}"]
        &&
    {\Exp\vB\vA{\vC'}}
        \arrow[d, "{\Lam{\va'} (\dash)_{\vC'}\tup{1_{\vC'}}{\va'}}"]
    \\
    {\Exp{\vB\vC}{\vA\vC}}
        \arrow[dr, "{\vB\vh \comp \dash}"]
        &&
    {\Exp{\vB\vC'}{\vA\vC'}}
        \arrow[dl, "{\dash \comp \vA\vh}"]
    \\
    &
    {\Exp{\vB\vC'}{\vA\vC}}
    \end{tikzcd}
\end{center}

The morphism  $\End{\vf : \vD \to \vF\vC \in \commaC\vD\vF} \Exp{\vB\vC}{\vA\vC} \dashrightarrow \End{\vf : \vD \to \vF\vC \in \commaC\vD\vF} \Exp\vB\vA\vC$ is from the following wedge:
\begin{center}
    \begin{tikzcd}
        &
        {\End{\vf : \vD \to \vF\vC \in \commaC\vD\vF} \Exp{\vB\vC}{\vA\vC}}
            \arrow[dl, "{(\Lam{\tup\vk\va} \proj{\vF\vk \comp \vf} (\dash)(\va))_{\hat\vC \in \mcatC}}"']
            \arrow[dr, "{(\Lam{\tup\vk\va} \proj{\vF\vk \comp \vg} (\dash)(\va))_{\hat\vC \in \mcatC}}"]
    \\
        {\Exp\vB\vA\vC}
        \arrow[rr, "{\Exp\vB\vA\vh = (\Lam{\tup\vk\va} (\dash)_{\hat\vC}\tup{\vk \comp \vh}\va)_{\hat\vC \in \mcatC}}"']
    &&
        {\Exp\vB\vA{\vC'}}
    \end{tikzcd}
\end{center}
    which commutes by:
\begin{tikzcd}
    {\vF{\hat\vC}}
        \arrow[rr, no head, Rightarrow, no head]
    &&
    {\vF{\hat\vC}}
    \\
    {\vF{\vC'}}
        \arrow[u, "{\vF\vk}"]
    &&
    \\
    {\vF\vC}
        \arrow[u, "{\vF\vh}"]
        \arrow[rr, "{\vF\vh}"]
    &&
    {\vF\vC'}
        \arrow[uu, "{\vF\vk}"]
    \\
    &
    {\vD}
        \arrow[ul, "\vf"]
        \arrow[ur, "\vg"']
\end{tikzcd}
\end{proof}
\noindent
That the end projections form a natural transformation and that the two induced morphisms are inverses is not hard to prove.
